%-*-latex-*-
%-*-latex-*-
\newcommand\COURSE{ciss245}
\newcommand\ASSESSMENT{q8101}
\newcommand\ASSESSMENTTYPE{Quiz}
\newcommand\POINTS{\textwhite{xxx/xxx}}

\input{myquizpreamble}
\input{yliow}
\input{\COURSE}
%-*-latex-*-

%-*-latex-*-

%-*-latex-*-
\renewcommand\TITLE{\ASSESSMENTTYPE \ \ASSESSMENT}

\renewcommand\EMAIL{}
\renewcommand\AUTHOR{eagraham1@cougars.ccis.edu}

\textwidth=6in
\begin{document}
\topmatterthree

%------------------------------------------------------------------------------
\nextq
Make the least possible modification to the given program so that the output is
\begin{console}[fontsize=\footnotesize]
03
\end{console}
Do \textit{not} modify \verb!main()!.
\\
\ANSWER
\begin{answercode}
#include <iostream>

class A
{
public:
    void m()
    {
        std::cout << '0';
    }
    virtual void n()
    {
        std::cout << '1';
    }
};

class B: public A
{
public:
    void m()
    {
        std::cout << '2';
    }
    void n()
    {
        std::cout << '3';
    }
};

int main()
{
    A * p = new B;
    p->m();
    p->n();
    return 0;
}
\end{answercode}

%------------------------------------------------------------------------------
\nextq
John wants to put 42 and 3.14 into a vector and print their values.
But the code below (obviously!) does not work.
Help him out by correcting all the mistakes.
(Hint: create a parent class \verb!Number! for both \verb!Int!
and \verb!Double! and use polymorphism.)
\\
\ANSWER
\begin{answercode}
#include <iostream>
#include <vector>

// Put a parent class here
class Number
{
  public:
  virtual void print() const;
}

class Int: public Number
{
public:
    Int(int v)
        : v_(v)
    {}
    void print() const
    {
        std::cout << v_ << '\n';
    }
    int v_;
};

class Double: public Number
{
public:
    Double(double v)
        : v_(v)
    {}
    void print() const
    {
        std::cout << v_ << '\n';
    }
    double v_;
};

int main()
{
    std::vector< Number > numbers;
    numbers.push_back(Int(42));
    numbers.push_back(Double(3.14));

    for (int i = 0; i < 2; ++i)
    {
        numbers[i].print();
    }

    return 0;
}
\end{answercode}

%------------------------------------------------------------------------------
%-*-latex-*-
\newpage
\input{instructions}
\end{document}

